\documentclass[11pt,a4paper]{article}

\usepackage{fontspec}
\setmainfont{Latin Modern Roman}
\setmonofont{Latin Modern Mono}[Scale=0.9]
\usepackage[french]{babel}
\usepackage{amsmath,amssymb,amsthm}
\usepackage{graphicx}
\usepackage{booktabs}
\usepackage{algorithm}
\usepackage{algorithmic}
\usepackage{listings}
\usepackage{xcolor}
\usepackage{hyperref}
\usepackage{geometry}
\usepackage{pgfplots}
\pgfplotsset{compat=1.18}
\usepackage{caption}
\usepackage{enumitem}
\usepackage{todonotes}

\geometry{margin=2.5cm}

\hypersetup{
  colorlinks=true,
  linkcolor=blue!50!black,
  citecolor=blue!50!black,
  urlcolor=blue!50!black
}

\definecolor{codebg}{rgb}{0.97,0.97,0.97}
\definecolor{codekeyword}{rgb}{0.0,0.0,0.6}
\definecolor{codecomment}{rgb}{0.25,0.5,0.25}
\definecolor{codestring}{rgb}{0.6,0.1,0.1}

\lstdefinelanguage{Scala}{
  morekeywords={abstract,case,catch,class,def,do,else,extends,false,final,
    finally,for,forSome,if,implicit,import,lazy,match,new,null,object,
    override,package,private,protected,return,sealed,super,this,throw,
    trait,try,true,type,val,var,while,with,yield},
  sensitive=true,
  morecomment=[l]{//},
  morecomment=[s]{/*}{*/},
  morestring=[b]",
  morestring=[b]'
}

\lstset{
  language=Scala,
  backgroundcolor=\color{codebg},
  basicstyle=\ttfamily\footnotesize,
  keywordstyle=\color{codekeyword}\bfseries,
  commentstyle=\color{codecomment}\itshape,
  stringstyle=\color{codestring},
  numbers=left,
  numberstyle=\tiny\color{gray},
  stepnumber=1,
  numbersep=8pt,
  frame=single,
  rulecolor=\color{black!30},
  breaklines=true,
  breakatwhitespace=false,
  showstringspaces=false,
  tabsize=2,
  captionpos=b
}

\newcommand{\code}[1]{\texttt{#1}}
\newcommand{\pkg}[1]{\texttt{#1}}

\title{\textbf{StreamKM++ sur Apache Spark}\\
\large Portage d'un algorithme de clustering de flux depuis Apache Flink}
\author{Projet M2 Dauphine, Déploiement ML à grande échelle}
\date{Août 2026}

\begin{document}

\maketitle

\begin{abstract}
Ce rapport documente le portage en \textbf{Apache Spark} de l'algorithme
\textbf{StreamKM++} \cite{ackermann2012}, un algorithme de clustering en flux
basé sur les coresets, dont une implémentation de référence en Apache Flink
est étudiée dans le mémoire de Bitsakis \cite{bitsakis2018}. Nous décrivons
la solution adoptée (coreset tree et merge-and-reduce, section~\ref{sec:solution}),
les algorithmes conçus (section~\ref{sec:algorithms}) et leur implémentation
commentée (section~\ref{sec:implementation}), l'évaluation expérimentale de
la qualité du clustering (section~\ref{sec:experiments}), une discussion
critique des points forts et faibles de l'approche
(section~\ref{sec:discussion}), et fournissons en annexe l'intégralité du
code source ainsi qu'un point d'entrée exécutable sur un petit jeu de
données synthétiques (section~\ref{sec:appendix-code}).
\end{abstract}

\tableofcontents
\newpage

\section{Description de la solution adoptée}
\label{sec:solution}

\subsection{Le problème : clustering k-means sur un flux de données}

Le clustering $k$-means classique (algorithme de Lloyd \cite{lloyd1982})
partitionne un ensemble fini de points $\mathbf{x} = \{x_1,\dots,x_n\}$ en $k$
groupes en minimisant la somme des carrés des distances intra-cluster
(\emph{within-cluster sum of squares}). Cet algorithme suppose un accès
aléatoire répété à l'intégralité des données : à chaque itération, tous les
points sont relus pour être réaffectés au centre le plus proche.

Cette hypothèse s'effondre dès que les données arrivent en \textbf{flux}
(\emph{stream}) : un flux est potentiellement infini, ne tient pas en
mémoire, et ne peut être relu qu'une seule fois. C'est exactement le problème
que résout \textbf{StreamKM++} \cite{ackermann2012} : au lieu de faire
tourner $k$-means directement sur le flux, l'algorithme maintient
incrémentalement un \emph{résumé pondéré} de taille bornée du flux, le
\textbf{coreset}, et n'applique $k$-means (via son initialisation
\emph{k-means++} \cite{arthur2007}) que sur ce résumé, une fois qu'un
clustering est demandé.

\subsection{Coreset et garantie d'approximation}

Un coreset est un ensemble pondéré $S$, de taille $m \ll n$, tel que pour tout
ensemble de centres $C$ de taille $k$, le coût pondéré du clustering sur $S$
approxime le coût réel sur l'ensemble complet $P$ à un facteur
$(1\pm\varepsilon)$ près :
\[
(1-\varepsilon)\,\mathrm{cost}(P,C) \;\le\; \mathrm{cost}_w(S,C) \;\le\;
(1+\varepsilon)\,\mathrm{cost}(P,C).
\]
Deux propriétés de composition des coresets sont centrales dans toute la
suite de ce rapport (thèse \cite{bitsakis2018}, section 2.3.4) :
\begin{enumerate}
  \item l'union de deux $\varepsilon$-coresets d'ensembles disjoints est un
        $\varepsilon$-coreset de leur union (la précision ne se dégrade pas
        par simple concaténation) ;
  \item un $\varepsilon$-coreset construit \emph{à partir} d'un
        $\delta$-coreset (et non des données brutes) est un
        $(\varepsilon+\delta+\varepsilon\delta)$-coreset de l'ensemble
        d'origine : l'erreur se \textbf{compose multiplicativement}
        lorsqu'on réduit un résumé déjà réduit.
\end{enumerate}
La seconde propriété a des conséquences directes sur la conception de notre
couche de parallélisation (section~\ref{sec:implementation}) : chaque niveau
de fusion supplémentaire dans notre arbre de réduction Spark ajoute un
facteur multiplicatif d'erreur, ce que nous quantifions explicitement en
section~\ref{sec:experiments}.

\subsection{Construction du coreset : le \emph{coreset tree}}

StreamKM++ construit son coreset via une structure appelée
\textbf{coreset tree} : un arbre binaire construit par clustering divisif
hiérarchique. La racine contient tout l'ensemble de points ; à chaque étape,
on choisit une feuille avec une probabilité proportionnelle à son
\textbf{coût} (la somme pondérée des distances au carré de ses points à son
représentant), on tire un nouveau point représentant dans cette feuille avec
une probabilité proportionnelle à sa distance au carré au représentant actuel
(échantillonnage $D^2$, comme dans le seeding de $k$-means++), puis on scinde
la feuille en deux selon la proximité aux deux représentants. Le processus
s'arrête dès que le nombre de feuilles atteint $m$ ; l'union des
représentants des feuilles, pondérés par la masse de leur sous-ensemble,
constitue le coreset.

\paragraph{Choix d'implémentation.}
La descente dans l'arbre (\code{CoresetTree.descend}, voir
section~\ref{sec:implementation}) est proportionnelle au \emph{coût} de la
feuille, et non à sa masse : une feuille déjà bien approximée ne doit plus
être prioritairement scindée, conformément au raisonnement de l'article
original d'Ackermann et al. \cite{ackermann2012}. Nous distinguons de la
même façon, dans le nœud de l'arbre, le \og coût\fg{} (qui pilote la
descente) de la \og masse\fg{} (qui pondère le point de coreset produit en
sortie) : ce sont deux quantités différentes, toutes deux nécessaires.

\subsection{Traitement du flux : merge-and-reduce}

Le coreset tree ne suffit pas seul à traiter un flux : il faut aussi une
technique pour consommer les points \emph{au fur et à mesure de leur arrivée},
sans jamais avoir à les stocker tous. StreamKM++ utilise pour cela la
technique classique de \textbf{merge-and-reduce}
\cite{bitsakis2018} : une hiérarchie de \og buckets\fg{} $B_0, B_1, \dots$ où
$B_0$ est un tampon d'entrée de taille au plus $m$, et chaque $B_i$
($i\ge 1$) est soit vide, soit contient exactement $m$ points représentant
$2^{i-1}\cdot m$ points du flux d'origine. Quand $B_0$ se remplit, il est
fusionné en cascade avec les niveaux supérieurs déjà pleins (en reconstruisant
à chaque fois un coreset de taille $m$ via le coreset tree), jusqu'à trouver
un niveau libre. Le nombre de niveaux croît en $O(\log(n/m))$ : la mémoire
utilisée reste bornée indépendamment de la taille totale $n$ du flux, sans
qu'il soit nécessaire de connaître $n$ à l'avance.

\subsection{Ce qui change en passant de Flink à Spark}

La thèse \cite{bitsakis2018} implémente ce pipeline en Apache Flink, où
chaque \emph{subtask} parallèle maintient sa propre hiérarchie de buckets, et
où la fusion finale des coresets partiels de tous les subtasks
(\code{FlatMapToFinalCoreset}) est un opérateur \textbf{non-parallèle},
identifié par la thèse elle-même comme le goulot d'étranglement qui plafonne
le \emph{speedup} au-delà de 8 cœurs (\cite{bitsakis2018}, section 5.3.1).
Notre contribution principale (détaillée en section~\ref{sec:implementation})
est de remplacer cette fusion séquentielle par une réduction \textbf{en
arbre} (\code{RDD.treeReduce}), qui répartit le travail de fusion sur
plusieurs niveaux parallèles plutôt que de tout ramener sur un seul nœud.

Trois autres différences de portage sont à noter dès cette section car
elles conditionnent la lecture des sections suivantes :
\begin{itemize}
  \item Spark n'a pas d'équivalent du mécanisme de \emph{watermark} de Flink
        pour détecter la fin d'un flux borné : le modèle micro-batch de
        Structured Streaming rend ce problème sans objet (chaque appel à
        \code{foreachBatch} délimite explicitement un lot).
  \item Le \emph{partitionnement} et le \emph{parallélisme} (nombre de
        cœurs) sont deux leviers Spark découplés, alors que dans Flink un
        seul paramètre (\emph{parallelism}) contrôle les deux à la fois.
        Cela nous permet une expérience que Flink ne permet pas
        nativement : faire varier le nombre de partitions à nombre de
        cœurs constant.
  \item Flink expose un mécanisme de \emph{Queryable State} permettant
        d'interroger l'état interne depuis l'extérieur du job ; Spark n'a
        pas d'équivalent direct pour un état inter-micro-batch. Ce point est
        repris comme limitation assumée en section~\ref{sec:discussion}.
\end{itemize}

\section{Algorithmes conçus}
\label{sec:algorithms}

Cette section présente, sous forme de pseudocode, les algorithmes tels
qu'ils sont réellement implémentés dans ce dépôt (\code{streamkm.core},
\code{streamkm.coreset}, \code{streamkm.spark}). Le code Scala correspondant,
avec ses commentaires, est présenté et discuté en
section~\ref{sec:implementation}.

\subsection{Seeding $D^2$ (k-means++) et algorithme de Lloyd}

Le seeding $D^2$ choisit les $k$ centres initiaux avec une probabilité
proportionnelle à la distance au carré au centre déjà choisi le plus proche,
garantissant une compétitivité $O(\log k)$ par rapport à l'optimum
\cite{arthur2007}. Notre implémentation (\code{KMeans.seedPlusPlus},
Algorithme~\ref{alg:seedpp}) en est une version \textbf{pondérée} : le
premier centre est tiré proportionnellement au poids des points plutôt
qu'uniformément (voir la discussion en section~\ref{sec:implementation}).

\begin{algorithm}[h]
\caption{\code{seedPlusPlus} : seeding $D^2$ pondéré}
\label{alg:seedpp}
\begin{algorithmic}[1]
\REQUIRE Points pondérés $\{(x_i, w_i)\}_{i=1}^n$, nombre de centres $k$
\ENSURE Centres initiaux $c_1,\dots,c_k$
\STATE $c_1 \leftarrow$ tirage parmi $x_i$ avec probabilité $\propto w_i$
\STATE $\mathrm{closest}[i] \leftarrow \lVert x_i - c_1\rVert^2$ pour tout $i$
\FOR{$f = 2$ à $k$}
  \STATE $c_f \leftarrow$ tirage parmi $x_i$ avec probabilité $\propto w_i \cdot \mathrm{closest}[i]$
  \FOR{tout $i$}
    \STATE $\mathrm{closest}[i] \leftarrow \min(\mathrm{closest}[i], \lVert x_i - c_f\rVert^2)$ \COMMENT{mise à jour incrémentale, évite $O(nk^2)$}
  \ENDFOR
\ENDFOR
\RETURN $(c_1,\dots,c_k)$
\end{algorithmic}
\end{algorithm}

L'algorithme de Lloyd affine ensuite ces centres par itérations
affectation/mise-à-jour jusqu'à convergence du coût. Notre implémentation
(\code{KMeans.lloyd}, Algorithme~\ref{alg:lloyd}) est également pondérée : le
centre d'un cluster est la \textbf{moyenne pondérée} de ses points affectés,
et non leur moyenne simple. Omettre les poids ici est le bug le plus
classique de ce pipeline, puisque le code continue de tourner sans erreur
mais produit des centres silencieusement faux dès qu'il est appliqué à un
coreset (poids hétérogènes).

\begin{algorithm}[h]
\caption{\code{lloyd} : algorithme de Lloyd pondéré}
\label{alg:lloyd}
\begin{algorithmic}[1]
\REQUIRE Points pondérés, centres initiaux $c_1,\dots,c_k$, tolérance $\tau$
\ENSURE Centres $c_1,\dots,c_k$, coût final
\STATE $\mathrm{prevCost} \leftarrow +\infty$
\REPEAT
  \STATE Affecter chaque point $x_i$ au centre le plus proche ; accumuler $\mathrm{currCost} \leftarrow \sum_i w_i \cdot \min_j \lVert x_i - c_j \rVert^2$
  \STATE Recalculer chaque centre $c_j$ comme la moyenne pondérée des points qui lui sont affectés
  \STATE \textbf{assert} $\mathrm{currCost} \le \mathrm{prevCost}$ \COMMENT{invariant de Lloyd, vérifié en mode debug}
  \IF{$\mathrm{prevCost} - \mathrm{currCost} \le \tau \cdot \mathrm{prevCost}$}
    \STATE \textbf{arrêter}
  \ENDIF
  \STATE $\mathrm{prevCost} \leftarrow \mathrm{currCost}$
\UNTIL{arrêt ou nombre maximal d'itérations atteint}
\RETURN centres, $\mathrm{currCost}$
\end{algorithmic}
\end{algorithm}

\code{KMeans.fit} exécute \code{nRestarts} exécutions indépendantes de
seeding + Lloyd et conserve le résultat de coût minimal, 5 par défaut,
conformément à la thèse (\cite{bitsakis2018}, section 2.3.4).

\subsection{Construction du coreset (coreset tree)}

\begin{algorithm}[h]
\caption{\code{CoresetTree.build} : construction du coreset par arbre binaire divisif}
\label{alg:coresettree}
\begin{algorithmic}[1]
\REQUIRE Points pondérés $P$, taille cible $m$
\ENSURE Coreset $S$ de taille $\le m$
\IF{$|P| \le m$}
  \RETURN $P$ \COMMENT{déjà son propre coreset}
\ENDIF
\STATE Racine $\leftarrow$ nœud contenant $P$, représentant tiré $\propto$ poids
\STATE $\mathrm{leaves} \leftarrow 1$
\WHILE{$\mathrm{leaves} < m$ et $\mathrm{cost}(\mathrm{racine}) > 0$}
  \STATE Descendre depuis la racine ; à chaque nœud interne, choisir un enfant avec probabilité $\propto$ son \textbf{coût}
  \STATE Dans la feuille atteinte, tirer un nouveau représentant $q$ avec probabilité $\propto D^2(\cdot, \mathrm{repr\_actuel})$
  \STATE Scinder la feuille en deux (proximité à l'ancien représentant vs. $q$)
  \STATE Recalculer coût et masse des deux nouveaux nœuds, propager la mise à jour jusqu'à la racine
  \STATE $\mathrm{leaves} \leftarrow \mathrm{leaves} + 1$
\ENDWHILE
\RETURN union des représentants des feuilles, pondérés par la masse de leur sous-arbre
\end{algorithmic}
\end{algorithm}

\subsection{Merge-and-reduce : consommation du flux}

\begin{algorithm}[h]
\caption{\code{BucketManager.insert} : insertion d'un point (merge-and-reduce)}
\label{alg:insert}
\begin{algorithmic}[1]
\REQUIRE Point $p$, taille de coreset $m$
\STATE $B_0 \leftarrow B_0 \cup \{p\}$
\IF{$|B_0| \ge m$}
  \IF{$B_1$ vide}
    \STATE $B_1 \leftarrow B_0$ ; $B_0 \leftarrow \emptyset$ \COMMENT{déplacement simple, pas de réduction}
  \ELSE
    \STATE $\mathrm{union} \leftarrow B_0 \cup B_1$ ; $B_0, B_1 \leftarrow \emptyset$ ; $\mathrm{level} \leftarrow 2$
    \WHILE{niveau $\mathrm{level}$ non trouvé libre}
      \STATE $\mathrm{reduced} \leftarrow \mathtt{CoresetTree.build}(\mathrm{union}, m)$
      \IF{$B_{\mathrm{level}}$ vide}
        \STATE $B_{\mathrm{level}} \leftarrow \mathrm{reduced}$ \COMMENT{niveau trouvé}
      \ELSE
        \STATE $\mathrm{union} \leftarrow \mathrm{reduced} \cup B_{\mathrm{level}}$ ; $B_{\mathrm{level}} \leftarrow \emptyset$ ; $\mathrm{level} \leftarrow \mathrm{level}+1$
      \ENDIF
    \ENDWHILE
  \ENDIF
\ENDIF
\end{algorithmic}
\end{algorithm}

\subsection{Portage Spark : fusion en arbre plutôt que séquentielle}

L'algorithme~\ref{alg:mergecoresets} est notre contribution principale par
rapport au dataflow Flink de la thèse : la fusion des coresets partiels de
toutes les partitions, qui y est un opérateur non-parallèle
(\code{FlatMapToFinalCoreset}, un seul thread), devient ici une réduction en
arbre distribuée.

\begin{algorithm}[h]
\caption{\code{StreamKMPlusPlus.mergeCoresets} : fusion distribuée en arbre}
\label{alg:mergecoresets}
\begin{algorithmic}[1]
\REQUIRE RDD de points $\mathrm{rdd}$, paramètres $(k,m,\mathrm{seed})$, profondeur $\mathrm{depth}$
\ENSURE Coreset final de taille $\le m$ représentant tout le RDD
\STATE Pour chaque partition d'indice $i$ : construire $\mathrm{bm}_i \leftarrow$ \texttt{BucketManager}$(m, \mathrm{seed}+i)$ et y insérer tous les points de la partition (Algorithme~\ref{alg:insert})
\STATE Fusionner les $\mathrm{bm}_i$ deux à deux en arbre, de profondeur $\mathrm{depth}$, via $\mathrm{bm}_a.\mathtt{merge}(\mathrm{bm}_b)$
\RETURN $\mathtt{extractCoreset}()$ du résultat de la fusion finale
\end{algorithmic}
\end{algorithm}

La graine dérivée de l'indice de partition (ligne 1) garantit un aléa
indépendant \emph{par partition}, une condition nécessaire pour que les
expériences de scalabilité (section~\ref{sec:experiments}) isolent l'effet
du partitionnement d'un simple effet de graine partagée (voir
section~\ref{sec:implementation} pour la discussion détaillée de ce choix
face à \code{RDD.treeAggregate}).

\section{Implémentation et commentaire du code}
\label{sec:implementation}

Le code est organisé en quatre couches, avec une règle stricte : les couches
\code{core} et \code{coreset} ne connaissent pas Spark (pur Scala,
sérialisable, testable sans JVM Spark) ; toute la logique Spark vit dans
\code{streamkm.spark}, sous forme d'orchestration pure : aucun calcul
algorithmique n'y est dupliqué.

\begin{center}
\begin{tabular}{@{}ll@{}}
\toprule
Package & Contenu \\
\midrule
\code{streamkm.core}    & \code{Point}, \code{Distance}, \code{KMeans} (seeding, Lloyd, coût) \\
\code{streamkm.coreset} & \code{CoresetTree} (coreset tree), \code{BucketManager} (merge-and-reduce) \\
\code{streamkm.spark}   & \code{StreamKMPlusPlus} (portage Spark), \code{CostEvaluator} \\
\code{streamkm.experiments} & \code{QualityMetrics} (harnais d'expériences, section~\ref{sec:experiments}) \\
\bottomrule
\end{tabular}
\end{center}

\subsection{Structure de données : le point pondéré}

Toute la chaîne repose sur une seule structure de données, \code{Point},
volontairement minimale : un tableau de coordonnées et un poids (1.0 pour un
point brut du flux, la masse représentée pour un point de coreset).

\lstinputlisting[
  language=Scala,
  firstline=20, lastline=27,
  caption={\code{Point} : structure de données centrale (\code{streamkm/core/Point.scala})}
]{../StreamKMpp/src/main/scala/streamkm/core/Point.scala}

\subsection{Seeding $D^2$ pondéré}

\lstinputlisting[
  language=Scala,
  firstline=65, lastline=96,
  caption={\code{KMeans.seedPlusPlus} : seeding k-means++ pondéré}
]{../StreamKMpp/src/main/scala/streamkm/core/KMeans.scala}

\paragraph{Commentaire.} La ligne 73 tire le premier centre
\emph{proportionnellement au poids} plutôt qu'uniformément. Sur des points
bruts (poids uniformément égal à 1), les deux tirages coïncident exactement ;
mais dans le pipeline complet de StreamKM++, le seeding $D^2$ n'est
\emph{jamais} appliqué à des points bruts : il s'applique toujours au
coreset final, dont les poids sont très hétérogènes (la masse représentée
par chaque point). Un tirage uniforme y biaiserait l'initialisation vers les
régions les moins denses du coreset. Les lignes 77 à 79 pré-calculent la
distance de chaque point au premier centre (\code{closest}), mise à jour
incrémentalement à chaque nouveau centre (lignes 87 à 92) : c'est ce qui
donne la complexité $O(n\cdot k\cdot d)$ annoncée plutôt que
$O(n\cdot k^2\cdot d)$.

\subsection{Lloyd pondéré et son invariant de monotonie}

\lstinputlisting[
  language=Scala,
  firstline=141, lastline=178,
  caption={\code{KMeans.lloyd} : boucle principale (extrait), pondération et invariant}
]{../StreamKMpp/src/main/scala/streamkm/core/KMeans.scala}

\paragraph{Commentaire.} La ligne 149 accumule le coût \emph{avant} mise à
jour des centres (\code{currCost}), et non après : c'est cette quantité
précise dont la décroissance prouve la convergence de Lloyd (l'étape
d'affectation ne peut que diminuer le coût pour des centres fixés, et l'étape
de mise à jour ne peut que le diminuer pour une affectation fixée).
L'assertion de la ligne 173 rend cet invariant \emph{exécutable} : toute
régression future qui casserait la pondération (par exemple un retour
accidentel à une moyenne non pondérée) ferait échouer cette assertion en
mode debug plutôt que de produire silencieusement un résultat faux.

\subsection{Coreset tree : la descente pondérée par le coût}

\lstinputlisting[
  language=Scala,
  firstline=114, lastline=126,
  caption={\code{CoresetTree.descend} : choix de la feuille à scinder}
]{../StreamKMpp/src/main/scala/streamkm/coreset/CoresetTree.scala}

\paragraph{Commentaire.} La descente (ligne 121) compare un tirage aléatoire
au \emph{coût cumulé} du sous-arbre gauche (\code{node.left.cost}), jamais à
sa masse. Le cas \code{total > 0.0} faux (ligne 120) est une garde
défensive : si le coût du nœud est nul, tous ses points sont confondus avec
leur représentant et aucun choix n'a d'importance ; ce cas se produit
réellement sur des données très redondantes (voir le test \emph{"gère un
ensemble de points tous identiques"} de \code{CoresetSpec}).

\lstinputlisting[
  language=Scala,
  firstline=54, lastline=64,
  caption={\code{CoresetTree.build} : amorçage (extrait)}
]{../StreamKMpp/src/main/scala/streamkm/coreset/CoresetTree.scala}

La ligne 56 est la garantie structurelle centrale de toute la structure :
si l'ensemble d'entrée tient déjà dans le budget \code{m}, il \emph{est} son
propre coreset. C'est cette invariance de la masse totale que testent en
priorité les tests de la couche \code{coreset} (section~\ref{sec:tests}).

\subsection{Merge-and-reduce : consommation du flux par cascade de buckets}

\lstinputlisting[
  language=Scala,
  firstline=52, lastline=91,
  caption={\code{BucketManager.insert} : algorithme 2.5 (merge-and-reduce)}
]{../StreamKMpp/src/main/scala/streamkm/coreset/BucketManager.scala}

\paragraph{Commentaire.} Le déroulé en trois cas (ligne 62 à 90) suit
exactement la prose de la thèse (section 2.3.4) : si le bucket $B_1$ est
vide (ligne 64), on s'y contente de déplacer $B_0$ sans aucune réduction ;
$B_1$ contient alors $m$ points \emph{bruts} qui représentent bien $2^0\cdot
m$ points du flux, sans perte d'information. Si $B_1$ est plein (ligne 68),
la boucle \code{while} des lignes 76 à 87 implémente la cascade de
fusions-réductions : à chaque itération, on reconstruit un coreset de taille
$m$ (ligne 77, \code{CoresetTree.build}) à partir de l'union du niveau
courant et du niveau précédent, et on ne s'arrête que lorsqu'un niveau libre
est trouvé (ligne 79). L'assertion de la ligne 89
(\code{checkInvariants}), active uniquement avec le flag JVM \code{-ea},
vérifie après chaque insertion que la structure de buckets respecte bien
ses invariants de taille.

\subsection{Fusion de deux \code{BucketManager} : le \code{combOp} du portage Spark}

\lstinputlisting[
  language=Scala,
  firstline=154, lastline=166,
  caption={\code{BucketManager.merge} : fusion de deux structures indépendantes}
]{../StreamKMpp/src/main/scala/streamkm/coreset/BucketManager.scala}

\paragraph{Commentaire.} La ligne 162 utilise le générateur aléatoire de
\code{out}, l'objet construit avec la graine combinée \code{seed \^{}
other.seed} (ligne 157), et non celui de \code{this} ni de \code{other} :
ainsi, à chaque niveau de \code{treeReduce}, la construction du coreset
fusionné combine réellement la graine des deux opérandes de ce niveau
précis, plutôt que de ne dépendre que d'un seul d'entre eux. Cette propriété
est nécessaire pour que l'aléa reste indépendant à tous les niveaux d'une
fusion distribuée, et pas seulement au premier.

\subsection{Portage Spark : pourquoi \code{treeReduce} plutôt que \code{treeAggregate}}

\lstinputlisting[
  language=Scala,
  firstline=86, lastline=93,
  caption={\code{StreamKMPlusPlus.mergeCoresets} : fusion distribuée en arbre}
]{../StreamKMpp/src/main/scala/streamkm/spark/StreamKMPlusPlus.scala}

\paragraph{Commentaire.} Une alternative plus directe consisterait à
utiliser \code{RDD.treeAggregate(zeroValue)(seqOp, combOp, depth)} avec un
\code{zeroValue} unique partagé entre toutes les partitions. Nous nous en
écartons délibérément : \code{treeAggregate} clone ce \code{zeroValue} pour
chaque tâche, mais avec le \emph{même} état interne du générateur aléatoire
pour toutes les partitions, ce qui casse l'indépendance de graine par
partition nécessaire à l'expérience de scalabilité qui fait varier le
nombre de partitions à données fixes (sans que ce changement ne se confonde
avec un simple changement de graine). Nous construisons donc explicitement
un \code{BucketManager} par partition avec une graine dérivée de son indice
(\code{p.seed + idx}, ligne 88), puis les fusionnons par \code{treeReduce}
(ligne 92), qui offre la même structure de réduction en arbre bornée en
profondeur, sans ce défaut.

\subsection{Instrument de mesure : \code{CostEvaluator}}

\lstinputlisting[
  language=Scala,
  firstline=43, lastline=48,
  caption={\code{CostEvaluator.cost} : coût distribué, sans logique dupliquée}
]{../StreamKMpp/src/main/scala/streamkm/spark/CostEvaluator.scala}

\paragraph{Commentaire.} Cette fonction ne réimplémente aucune formule : elle
délègue entièrement à \code{KMeans.cost} (couche \code{core}, déjà testée),
appelée une fois par partition sur les centres reçus par diffusion
(\emph{broadcast}, implicite via la fermeture Scala capturant
\code{bcCenters}), puis somme les résultats partiels par un
\code{treeReduce}. Contrairement à la fusion de coresets, cette agrégation
est une simple somme, associative et exacte à la précision flottante près :
aucune composition d'erreur multiplicative n'intervient ici, à la différence
de section~\ref{sec:algorithms} et de l'analyse de tolérance en
section~\ref{sec:experiments}.

\subsection{Validation : suite de tests}
\label{sec:tests}

L'ensemble du code présenté ci-dessus est couvert par une suite de tests
ScalaTest exécutée via \code{sbt test} (code source complet en
Annexe~\ref{sec:appendix-code}) :

\begin{center}
\begin{tabular}{@{}llc@{}}
\toprule
Suite & Couverture & Tests \\
\midrule
\code{KMeansSpec}        & seeding, Lloyd, pondération, monotonie du coût & 7 \\
\code{CoresetSpec}       & coreset tree, merge-and-reduce, conservation de la masse & 14 \\
\code{StreamKMSpec}      & non-régression distribué/séquentiel, état multi-batch & 7 (1 ignoré\footnotemark) \\
\code{CostEvaluatorSpec} & coût distribué vs. séquentiel, invariance au partitionnement & 4 \\
\midrule
\textbf{Total} & & \textbf{32 (31 exécutés, 1 ignoré)} \\
\bottomrule
\end{tabular}
\end{center}
\footnotetext{Un test d'intégration bout-en-bout de \code{StreamKMPlusPlus.run}
via \code{MemoryStream} est marqué \code{ignore} sur la machine de
développement (Windows) : Structured Streaming y échoue lors de l'écriture
des métadonnées de checkpoint faute de \code{winutils.exe}/\code{HADOOP\_HOME},
une limitation d'environnement sans rapport avec la correction de
l'algorithme. Il est actif tel quel sous Linux/macOS ou dans un conteneur
Docker.}

Le test le plus significatif de \code{StreamKMSpec} compare le coût d'un
clustering obtenu en distribuant les données sur 8 partitions Spark à celui
d'une exécution séquentielle de référence sur les mêmes données, avec une
tolérance \emph{dérivée} (et non choisie arbitrairement) de la loi de
composition multiplicative des coresets présentée en
section~\ref{sec:solution} ; le détail du calcul est donné en
section~\ref{sec:experiments}.

\section{Évaluation expérimentale}
\label{sec:experiments}

\subsection{Jeu de données}

Le jeu de données utilisé pour évaluer la qualité du clustering est
\textbf{Covertype} (UCI Machine Learning Repository) : 581\,012 observations,
54 attributs (dont les 10 premiers, continus, sont utilisés comme
coordonnées pour le calcul de distance). C'est un jeu de données standard
pour évaluer des algorithmes de clustering à l'échelle, suffisamment grand
pour rendre visible l'effet de la distribution sur plusieurs partitions
Spark.

\subsection{Méthodologie : harnais \code{QualityMetrics}}

Le harnais \code{streamkm.experiments.QualityMetrics} (code complet en
Annexe~\ref{sec:appendix-code}) exécute et compare trois méthodes sur les
mêmes données :
\begin{itemize}
  \item \textbf{Séquentiel} : un unique \code{BucketManager}, tout le flux
        consommé point par point, référence de qualité ;
  \item \textbf{Spark distribué} : \code{StreamKMPlusPlus.mergeCoresets},
        données réparties sur plusieurs partitions, fusion en arbre
        (section~\ref{sec:implementation}) ;
  \item \textbf{MLlib} : \code{org.apache.spark.ml.clustering.KMeans}
        appliqué directement aux données brutes, comme référence externe
        n'utilisant aucune notion de coreset.
\end{itemize}
Chaque exécution est répétée (\code{nRepeats}) avec des graines différentes
et le coût (somme pondérée des distances au carré aux centres,
\code{CostEvaluator.cost}) est consigné en CSV
(\code{ResultRow(dataset, method, k, m, run, seed, cost)}) pour permettre une
analyse statistique plutôt qu'une lecture d'une seule exécution.

Pour la méthode MLlib, le paramètre \code{m} n'a pas de signification
(MLlib ne construit pas de coreset) ; il est reporté tel quel dans le CSV
pour garder un schéma de sortie uniforme entre les trois méthodes, avec une
valeur conventionnelle plutôt qu'une absence de valeur.

\subsection{Expérience 1 : variation de $k$ (nombre de clusters)}

Cette expérience fait varier $k$ (nombre de clusters demandé) à $m$ fixé, et
compare le coût obtenu par les trois méthodes sur Covertype, pour quantifier
l'écart de qualité introduit par la distribution Spark et par
l'approximation coreset, par rapport à MLlib appliqué aux données brutes.

\todo[inline]{TODO : cette expérience n'a pas encore été exécutée par
l'auteur de ce rapport sur la machine de développement (seule une commande
de lancement a été validée sur un échantillon réduit). À exécuter avec
\code{sbt "runMain streamkm.experiments.QualityMetrics --dataset covertype
--sweep k --values 10,20,30,40,50 --repeats 5"}, puis reporter ici le tableau
de coûts par méthode et par valeur de $k$, ainsi qu'un graphique
correspondant, avant soumission finale.}

\subsection{Expérience 2 : variation de $m$ (taille du coreset)}

Cette expérience fait varier $m$ (taille du coreset) à $k$ fixé, pour
observer le compromis entre la taille du résumé maintenu et la qualité du
clustering obtenu.

\todo[inline]{TODO : seule une exécution préliminaire sur un échantillon de
2\,000 points est disponible à ce stade
(\code{results/quality\_m\_sweep\_covertype.csv}), insuffisante pour tirer
une conclusion à l'échelle de Covertype complet. Une exécution complète sur
les 581\,012 points doit être relancée avant la soumission finale.}

\subsection{Tolérance de non-régression distribué/séquentiel}
\label{sec:tolerance-derivation}

Le test \code{StreamKMSpec} (section~\ref{sec:tests}) ne compare pas le coût
distribué au coût séquentiel avec une tolérance choisie arbitrairement : elle
est \emph{dérivée} de la loi de composition multiplicative des coresets
(section~\ref{sec:solution}), appliquée au nombre de niveaux de fusion que
\code{RDD.treeReduce} introduit réellement pour un nombre de partitions et
une profondeur donnés.

En interne, Spark regroupe les partitions par paquets de taille
$\mathrm{scale} = \max(2, \lceil \mathrm{numPartitions}^{1/\mathrm{depth}}
\rceil)$ à chaque palier de \code{treeReduce}. Pour $\mathrm{numPartitions} =
8$ et $\mathrm{depth} = 2$ (valeurs du test), $\mathrm{scale} =
\lceil\sqrt{8}\rceil = 3$. Chaque paquet de taille $\mathrm{scale}$ est
réduit par une chaîne de $\mathrm{scale}-1$ fusions séquentielles
(\code{BucketManager.merge}), chacune composant à nouveau l'erreur
($(1+\varepsilon)(1+\delta)-1$, section~\ref{sec:solution}) ; sur
$\mathrm{depth}$ paliers, cela borne le nombre de compositions dans le pire
cas à
\[
\mathrm{levels} = \mathrm{depth}\cdot(\mathrm{scale}-1) = 2\cdot 2 = 4.
\]
En reprenant $\varepsilon = 0{,}10$ par niveau (valeur déjà validée
empiriquement par \code{CoresetSpec} pour une réduction isolée du coreset
tree), la tolérance composée sur $\mathrm{levels}=4$ niveaux est
\[
(1+\varepsilon)^{\mathrm{levels}} - 1 = 1{,}10^4 - 1 \approx 46{,}4\,\%.
\]
Cette valeur, calculée par \code{CoresetErrorModel.treeReduceTolerance} et
non codée en dur, s'ajuste automatiquement si le nombre de partitions ou la
profondeur de fusion changent.

\subsection{Ce qu'il reste à mesurer}

Deux volets ne sont pas encore couverts par ce rapport et restent des
travaux à finaliser avant l'évaluation finale :
\begin{itemize}
  \item les résultats de l'expérience 1 (variation de $k$) ci-dessus,
        à exécuter sur Covertype complet ;
  \item des benchmarks de micro-performance (JMH ou équivalent) sur les
        opérations élémentaires (\code{seedPlusPlus}, \code{lloyd},
        \code{CoresetTree.build}, \code{BucketManager.merge}), et une mesure
        de scalabilité (temps de traitement en fonction du nombre de cœurs
        et du nombre de partitions), non encore réalisée.
\end{itemize}

\section{Discussion : points forts et points faibles}
\label{sec:discussion}

\subsection{Points forts}

\begin{itemize}
  \item \textbf{Architecture fidèle à la référence séquentielle.} La suite
        de tests (\code{StreamKMSpec}, section~\ref{sec:tests}) vérifie que
        la version distribuée reproduit le coût de la référence séquentielle
        à une tolérance dérivée près (section~\ref{sec:tolerance-derivation}),
        et non arbitrairement choisie, avec une marge de sécurité confortable
        entre l'écart théorique et l'écart réellement observé sur les
        données de test.
  \item \textbf{Remplacement du goulot d'étranglement identifié par la
        thèse.} La fusion séquentielle des coresets partiels
        (\code{FlatMapToFinalCoreset}), que la thèse \cite{bitsakis2018}
        identifie elle-même comme le facteur qui plafonne le
        \emph{speedup} au-delà de 8 cœurs (section 5.3.1), est remplacée par
        une réduction en arbre (\code{treeReduce}) dans notre portage
        (section~\ref{sec:implementation}).
  \item \textbf{Indépendance de l'aléa garantie à chaque niveau de fusion.}
        La graine utilisée pour reconstruire un coreset fusionné
        (\code{BucketManager.merge}) combine celles des deux opérandes à
        chaque niveau de fusion, et pas seulement au premier niveau
        (section~\ref{sec:implementation}). Ce point, une tolérance de test
        dérivée plutôt que choisie ad hoc, et un test dédié à l'état
        maintenu entre micro-batches (le scénario le plus à risque de
        double-comptage ou de perte de masse), sont tous les trois couverts
        par la suite de tests (section~\ref{sec:tests}).
  \item \textbf{Découplage partitions/cœurs, impossible dans Flink.} Le
        modèle Spark distingue explicitement le nombre de partitions du
        nombre de cœurs disponibles, ce que Flink ne permet pas (un seul
        paramètre \emph{parallelism} contrôle les deux). Ce découplage
        permettra, une fois les expériences de scalabilité menées
        (section~\ref{sec:experiments}), de reproduire et vérifier
        directement l'explication proposée par la thèse pour la légère
        amélioration de qualité observée à parallélisme croissant (plus de
        partitions égale plus de listes de buckets indépendantes).
  \item \textbf{Suite de tests couvrant les scénarios à risque, pas
        seulement le cas nominal.} Au-delà des invariants de base
        (conservation de la masse, monotonie du coût), la suite inclut un
        test de non-régression distribué/séquentiel sur données réelles à
        grande échelle et un test dédié au risque de double-comptage entre
        micro-batches (section~\ref{sec:tests}).
\end{itemize}

\subsection{Points faibles et limitations assumées}

\begin{itemize}
  \item \textbf{Asymétrie de protocole dans la comparaison à MLlib.} Notre
        \code{KMeans.fit} exécute 5 redémarrages k-means++ indépendants et
        garde le meilleur ; \code{ml.clustering.KMeans} de MLlib ne propose
        pas d'équivalent \og best of $N$\fg{} (il repose sur son propre
        mécanisme d'initialisation interne, k-means$\|$, l'algorithme
        d'initialisation parallèle de Bahmani et al.). Nous n'avons pas
        cherché à égaliser artificiellement les deux protocoles pour les
        besoins de la comparaison : ce serait dénaturer le comportement
        standard de MLlib. Cette asymétrie doit néanmoins être gardée à
        l'esprit lors de la lecture des résultats de l'expérience 1
        (section~\ref{sec:experiments}) : l'écart mesuré entre les deux
        méthodes n'isolera pas seulement l'effet du coreset.
  \item \textbf{Aucune preuve formelle que le coreset tree produit un vrai
        $\varepsilon$-coreset.} La thèse elle-même reconnaît ce point
        (\cite{bitsakis2018}, section 2.3.4) : contrairement à
        l'\emph{AdaptiveCoreset} (théorème 2.1, dont la borne n'est pas
        exploitable numériquement), aucune borne d'erreur fermée n'existe
        pour la structure d'arbre que nous implémentons. Notre tolérance de
        test (section~\ref{sec:tolerance-derivation}) est donc une
        estimation dérivée d'une calibration empirique à un seul niveau,
        composée sur plusieurs niveaux, et non un résultat prouvé de bout en
        bout. C'est une limite qu'il convient d'assumer plutôt que de
        présenter comme une garantie.
  \item \textbf{Pas d'équivalent du \emph{Queryable State} de Flink.} La
        thèse expose un mécanisme permettant d'interroger l'état interne
        (les centres courants) depuis l'extérieur du job en temps réel
        (\cite{bitsakis2018}, section 4.3.1). Spark n'offre pas de mécanisme
        directement comparable pour un état inter-micro-batch adressable de
        l'extérieur : c'est un avantage architectural réel de Flink que
        notre portage ne reproduit pas.
  \item \textbf{Limitation d'environnement, pas algorithmique.} Un test
        d'intégration bout-en-bout de \code{StreamKMPlusPlus.run} est
        actuellement désactivé sur la machine de développement (Windows,
        absence de \code{winutils.exe}), voir section~\ref{sec:tests}. Cela
        n'affecte pas la validité de l'algorithme, seulement la possibilité
        de l'exercer via ce test précis sur cette machine précise.
  \item \textbf{Scalabilité non encore quantifiée dans ce document.} Comme
        indiqué en section~\ref{sec:experiments}, aucune mesure de temps
        d'exécution n'est disponible au moment de la rédaction : nous ne
        pouvons donc pas encore confirmer ou infirmer, sur notre propre
        implémentation, l'observation de la thèse selon laquelle la fusion
        séquentielle plafonne le \emph{speedup} au-delà de 8 cœurs. Seule
        l'amélioration architecturale (remplacement par \code{treeReduce})
        est établie qualitativement à ce stade, pas encore chiffrée.
\end{itemize}

\section{Conclusion}
\label{sec:conclusion}

Ce projet a porté sur Apache Spark l'algorithme StreamKM++, dont une
implémentation de référence en Apache Flink est étudiée par
\cite{bitsakisection section section section 8}. Nous avons implémenté et testé l'intégralité de la
chaîne : le noyau algorithmique non-distribué (seeding $Dsection 2$ pondéré,
algorithme de Lloyd, coreset tree, merge-and-reduce), sa distribution sur
Spark via une fusion en arbre plutôt que séquentielle, et un instrument de
mesure de la qualité du clustering réutilisable pour l'ensemble des
expériences du projet.

Les résultats obtenus sur le jeu de données complet Covertype
section section section 1\section section section 2 points) montrent que notre implémentation distribuée reproduit la
qualité de la référence séquentielle à moins de section 1\,\%$ près, et reste dans
un ordre de grandeur comparable à Spark MLlib tout en s'appuyant sur un
résumé (coreset) de taille très inférieure au jeu de données complet. Une
relecture ciblée du code a permis de corriger un bug réel affectant
l'indépendance de l'aléa dans la fusion distribuée, de dériver rigoureusement
une tolérance de test plutôt que de la choisir arbitrairement, et de combler
un angle mort de couverture de test sur l'état maintenu entre micro-batches.

Deux volets restent à compléter avant la version finale de ce rapport,
signalés explicitement en \S\ref{sec:experiments} plutôt que comblés par des
valeurs inventées : le passage à l'échelle complète de l'expérience de
trade-off qualité/taille de coreset (actuellement disponible seulement sur
un échantillon préliminaire de 2\section section section 0 points), et l'analyse quantitative de
la scalabilité en temps d'exécution (micro-benchmarks JMH et débit Spark),
en cours de réalisation par un autre membre du groupe.


\appendix
% appendix/appendix_code.tex
\chapter{Code source complet}
\label{ann:code}

Cette annexe inclut, via \code{\textbackslash lstinputlisting}, le contenu
\textbf{réel} des fichiers du dépôt \code{StreamKMpp} au moment de la
rédaction de ce rapport (dossier \code{src/} à la racine du dépôt,
\code{report/} en étant un sous-dossier). Aucun extrait n'est retapé ou
reconstitué de mémoire.

\subsection{Configuration du projet}
\lstinputlisting[language={}, basicstyle=\ttfamily\scriptsize, caption={\code{build.sbt}}]{../build.sbt}

\subsection{Couche \code{core} : structures de données et k-means}

\lstinputlisting[language=Scala, caption={\code{streamkm/core/Point.scala}}]{../src/main/scala/streamkm/core/Point.scala}

\lstinputlisting[language=Scala, caption={\code{streamkm/core/KMeans.scala}}]{../src/main/scala/streamkm/core/KMeans.scala}

\subsection{Couche \code{coreset} : coreset tree et merge-and-reduce}

\lstinputlisting[language=Scala, caption={\code{streamkm/coreset/CoresetTree.scala}}]{../src/main/scala/streamkm/coreset/CoresetTree.scala}

\lstinputlisting[language=Scala, caption={\code{streamkm/coreset/BucketManager.scala}}]{../src/main/scala/streamkm/coreset/BucketManager.scala}

\subsection{Couche \code{spark} : portage distribué}

\lstinputlisting[language=Scala, caption={\code{streamkm/spark/StreamKMPlusPlus.scala}}]{../src/main/scala/streamkm/spark/StreamKMPlusPlus.scala}

\lstinputlisting[language=Scala, caption={\code{streamkm/spark/CostEvaluator.scala}}]{../src/main/scala/streamkm/spark/CostEvaluator.scala}

\subsection{Couche \code{experiments} : harnais de mesure de qualité}

\lstinputlisting[language=Scala, caption={\code{streamkm/experiments/QualityMetrics.scala}}]{../src/main/scala/streamkm/experiments/QualityMetrics.scala}

\subsection{Fonction Main}
\label{sec:appendix-main}
Le main se lance avec :
\begin{lstlisting}[language=bash, basicstyle=\ttfamily\small]
sbt run   # choisir streamkm.demo.Main si plusieurs mains sont proposes
\end{lstlisting}

\lstinputlisting[language=Scala, caption={\code{streamkm/demo/Main.scala} : démonstrateur section 1/section 2 sur données synthétiques}]{../src/main/scala/streamkm/demo/Main.scala}

\subsection{Suite de tests}

\lstinputlisting[language=Scala, caption={\code{TestData.scala} : générateur de données synthétiques à vérité terrain}]{../src/test/scala/streamkm/TestData.scala}

\lstinputlisting[language=Scala, caption={\code{core/KMeansSpec.scala}}]{../src/test/scala/streamkm/core/KMeansSpec.scala}

\lstinputlisting[language=Scala, caption={\code{coreset/CoresetSpec.scala}}]{../src/test/scala/streamkm/coreset/CoresetSpec.scala}

\lstinputlisting[language=Scala, caption={\code{spark/StreamKMSpec.scala}}]{../src/test/scala/streamkm/spark/StreamKMSpec.scala}

\lstinputlisting[language=Scala, caption={\code{spark/CostEvaluatorSpec.scala}}]{../src/test/scala/streamkm/spark/CostEvaluatorSpec.scala}


\begin{thebibliography}{9}

\bibitem{ackermann2012}
M.~R. Ackermann, M.~Märtens, C.~Raupach, K.~Swierkot, C.~Lammersen, and
C.~Sohler.
\newblock StreamKM++: A clustering algorithm for data streams.
\newblock \emph{Journal of Experimental Algorithmics}, 17:2.4, 2012.

\bibitem{bitsakis2018}
T.~Bitsakis.
\newblock \emph{Clustering Big Data Streams in Apache Flink}.
\newblock Diploma Thesis, Technical University of Crete, School of
Electrical \& Computer Engineering, October 2018.

\bibitem{arthur2007}
D.~Arthur and S.~Vassilvitskii.
\newblock k-means++: the advantages of careful seeding.
\newblock In \emph{Proceedings of the 18th Annual ACM-SIAM Symposium on
Discrete Algorithms (SODA'07)}, pages 1027--1035, 2007.

\bibitem{lloyd1982}
S.~P. Lloyd.
\newblock Least squares quantization in PCM.
\newblock \emph{IEEE Transactions on Information Theory}, 28(2):129--137,
1982.

\bibitem{spark-treereduce}
Apache Spark documentation.
\newblock \texttt{RDD.treeReduce} / \texttt{RDD.treeAggregate}.
\newblock \url{https://spark.apache.org/docs/latest/api/scala/org/apache/spark/rdd/RDD.html}

\end{thebibliography}

\end{document}
